\documentclass[a4paper,12pt]{article}

% -------------------------------------------------
% Кодировка и язык
% -------------------------------------------------
\usepackage[T2A]{fontenc}
\usepackage[utf8]{inputenc}
\usepackage[russian]{babel}

% -------------------------------------------------
% Математика
% -------------------------------------------------
\usepackage{amsmath,amssymb,mathtools}
\usepackage{icomma}

% -------------------------------------------------
% Типографика
% -------------------------------------------------
\usepackage{helvet}
\renewcommand{\familydefault}{\sfdefault}
\usepackage[a4paper,top=2.0cm,bottom=2.1cm,left=2.1cm,right=2.1cm,headheight=15pt]{geometry}
\usepackage{microtype}
\usepackage{parskip}
\usepackage{setspace}
\onehalfspacing

% -------------------------------------------------
% Графика
% -------------------------------------------------
\usepackage{tikz}
\usetikzlibrary{calc}
\usepackage{pgfplots}
\pgfplotsset{compat=1.18}

% -------------------------------------------------
% Таблицы и списки
% -------------------------------------------------
\usepackage{tabularx,booktabs}
\usepackage{enumitem}

% -------------------------------------------------
% Служебное оформление
% -------------------------------------------------
\usepackage{xcolor}
\usepackage{fancyhdr}
\usepackage{hyperref}

% -------------------------------------------------
% Палитра
% -------------------------------------------------
\definecolor{accent}{HTML}{008080}
\definecolor{darkaccent}{HTML}{004D4D}
\definecolor{textgray}{HTML}{333333}
\definecolor{softgray}{HTML}{E8EEEE}
\definecolor{muted}{HTML}{657575}

% -------------------------------------------------
% Общие настройки
% -------------------------------------------------
\hypersetup{
    colorlinks=true,
    linkcolor=darkaccent,
    urlcolor=darkaccent
}

\setlength{\parindent}{0pt}
\setlength{\parskip}{0.65em}

\pagestyle{fancy}
\fancyhf{}
\fancyhead[L]{\small\color{muted}Графики функций с модулем}
\fancyhead[R]{\small\color{muted}ОГЭ}
\fancyfoot[C]{\small\color{muted}\thepage}
\renewcommand{\headrulewidth}{0.3pt}
\renewcommand{\headrule}{\hbox to\headwidth{\color{softgray}\leaders\hrule height \headrulewidth\hfill}}

\setlist[itemize]{
    leftmargin=1.5em,
    itemsep=0.2em,
    topsep=0.3em
}
\setlist[enumerate]{
    leftmargin=1.8em,
    itemsep=0.45em,
    topsep=0.3em
}

\newcommand{\lessonsection}[1]{%
    \vspace{0.6em}
    {\Large\bfseries\color{darkaccent}#1}\par
    \vspace{0.15em}
    {\color{softgray}\rule{\linewidth}{0.8pt}}
    \vspace{0.3em}
}

\newcommand{\important}[1]{%
    {\color{darkaccent}\bfseries #1}%
}

\newcommand{\R}{\mathbb{R}}

\begin{document}

% =================================================
% ТИТУЛЬНЫЙ ЛИСТ
% =================================================
\begin{titlepage}
\thispagestyle{empty}

\begin{center}

\vspace*{2.2cm}

{\Large\color{muted}ЧЕК-ЛИСТ}

\vspace{1.2cm}

{\Huge\bfseries\color{darkaccent}
Графики функций\\[0.15em]
с модулем}

\vspace{0.6cm}

{\Large
Кусочные функции, множество значений,\\
несколько модулей и алгебраические дроби}

\vspace{2.0cm}

{\large
София Кальней\\
ОГЭ по математике}

\vfill

{\large
\textbf{Лёвин Артём Александрович}\\[0.25em]
эксклюзивно для Софии Кальней}

\vspace{0.7cm}

{\large \today}

\vspace{2.2cm}

\end{center}

% Сдержанная кривая Безье
\begin{tikzpicture}[remember picture,overlay]
    \draw[
        color=accent!55,
        line width=1.3pt
    ]
    ([xshift=1.6cm,yshift=1.55cm]current page.south west)
    .. controls
    ([xshift=5.0cm,yshift=2.15cm]current page.south west)
    and
    ([xshift=-5.3cm,yshift=0.85cm]current page.south east)
    ..
    ([xshift=-1.6cm,yshift=1.50cm]current page.south east);
\end{tikzpicture}

\end{titlepage}

% =================================================
% ОСНОВНОЙ ЧЕК-ЛИСТ
% =================================================

\lessonsection{1. Главная идея: модуль превращает функцию в кусочную}

Для любого выражения \(A(x)\)
\[
|A(x)|=
\begin{cases}
 A(x), & A(x)\ge 0,\\[2mm]
 -A(x), & A(x)<0.
\end{cases}
\]

Поэтому при построении функции с модулем сначала нужно выяснить,
\important{при каких \(x\) подмодульное выражение положительно, отрицательно или равно нулю}.

\textbf{Базовый алгоритм}
\begin{enumerate}
    \item Найдите нули подмодульного выражения.
    \item Разбейте числовую прямую найденными точками на интервалы.
    \item Определите знак подмодульного выражения на каждом интервале.
    \item Раскройте модуль с соответствующим знаком.
    \item Получите кусочную формулу.
    \item Стройте каждую ветвь только на её собственном промежутке.
\end{enumerate}

\textbf{Важно.}
Граничную точку достаточно включить только в одну из соседних ветвей.
Например,
\[
x\le a,\qquad x>a,
\]
или
\[
x<a,\qquad x\ge a.
\]
Дважды включать одну и ту же точку в кусочную запись не требуется.

% -------------------------------------------------

\lessonsection{2. Пример: \(y=2-|x|\)}

Подмодульное выражение --- \(x\). Его нуль:
\[
x=0.
\]

Следовательно,
\[
y=2-|x|=
\begin{cases}
2+x, & x<0,\\[1mm]
2-x, & x\ge 0.
\end{cases}
\]

Обе ветви являются прямыми и встречаются в точке
\[
(0;2).
\]

\begin{center}
\begin{tikzpicture}
\begin{axis}[
    width=10.5cm,
    height=7.2cm,
    axis lines=middle,
    unit vector ratio=1 1,
    xmin=-4.5,
    xmax=4.5,
    ymin=-2.5,
    ymax=3.5,
    xtick={-4,-3,-2,-1,1,2,3,4},
    ytick={-2,-1,1,2,3},
    xlabel={$x$},
    ylabel={$y$},
    samples=250,
    clip=false,
    grid=both,
    major grid style={softgray},
    minor tick num=0
]
\addplot[
    color=accent,
    very thick,
    domain=-4.2:4.2
]
{2-abs(x)};
\addplot[
    only marks,
    mark=*,
    color=darkaccent
]
coordinates {(0,2)};
\node[anchor=south west,text=darkaccent]
at (axis cs:0.15,2.05) {$y=2-|x|$};
\end{axis}
\end{tikzpicture}
\end{center}

\textbf{Множество значений} --- все значения \(y\), которые достигает функция.

Из графика видно:
\[
E(y)=(-\infty;2].
\]

Значение \(2\) включается, поскольку вершина графика существует.

% -------------------------------------------------

\lessonsection{3. Как находить множество значений по графику}

Для проверки фиксируйте некоторое значение \(y=y_0\) и мысленно проводите
\important{горизонтальную прямую}.

\begin{itemize}
    \item Если прямая \(y=y_0\) пересекает график хотя бы один раз,
    то \(y_0\) входит в множество значений.
    \item Если пересечений нет, такое значение \(y_0\) не достигается.
\end{itemize}

Типичные варианты:
\[
(-\infty;a],\qquad
[a;+\infty),\qquad
\R.
\]

Для функции
\[
y=a-|x-b|
\]
точка \((b;a)\) является вершиной, поэтому
\[
E(y)=(-\infty;a].
\]

Для функции
\[
y=a+|x-b|
\]
получаем
\[
E(y)=[a;+\infty).
\]

% -------------------------------------------------

\lessonsection{4. Несколько модулей}

Если модулей несколько, нужно учесть нули \textbf{каждого}
подмодульного выражения.

Рассмотрим
\[
f(x)=|x-1|+|x+2|.
\]

Нули подмодульных выражений:
\[
x-1=0 \quad\Longrightarrow\quad x=1,
\]
\[
x+2=0 \quad\Longrightarrow\quad x=-2.
\]

Они разбивают числовую прямую на три промежутка:
\[
(-\infty;-2),\qquad [-2;1),\qquad [1;+\infty).
\]

\begin{center}
\begin{tikzpicture}[x=1.35cm,y=0.8cm]
    \draw[->,line width=0.9pt] (-4,0)--(4,0) node[right] {$x$};

    \draw (-2,0.11)--(-2,-0.11)
        node[below=5pt] {$-2$};
    \draw (1,0.11)--(1,-0.11)
        node[below=5pt] {$1$};

    \node[anchor=east] at (-4.0,1.35) {$x-1$};
    \node[anchor=east] at (-4.0,0.75) {$x+2$};

    \node at (-3,1.35) {$-$};
    \node at (-0.5,1.35) {$-$};
    \node at (2.4,1.35) {$+$};

    \node at (-3,0.75) {$-$};
    \node at (-0.5,0.75) {$+$};
    \node at (2.4,0.75) {$+$};
\end{tikzpicture}
\end{center}

Отсюда
\[
f(x)=
\begin{cases}
-(x-1)-(x+2), & x<-2,\\[1mm]
-(x-1)+(x+2), & -2\le x<1,\\[1mm]
(x-1)+(x+2), & x\ge 1.
\end{cases}
\]

После упрощения:
\[
\boxed{
f(x)=
\begin{cases}
-2x-1, & x<-2,\\[1mm]
3, & -2\le x<1,\\[1mm]
2x+1, & x\ge 1.
\end{cases}}
\]

\begin{center}
\begin{tikzpicture}
\begin{axis}[
    width=11cm,
    height=6.8cm,
    axis lines=middle,
    unit vector ratio=1 1,
    xmin=-4.5,
    xmax=4.5,
    ymin=-0.5,
    ymax=8,
    xlabel={$x$},
    ylabel={$y$},
    samples=250,
    clip=false,
    grid=both,
    major grid style={softgray}
]
\addplot[
    color=accent,
    very thick,
    domain=-4.2:4.2
]
{abs(x-1)+abs(x+2)};

\addplot[
    color=darkaccent,
    dashed,
    domain=-2:1
]
{3};

\node[anchor=south,text=darkaccent]
at (axis cs:-0.5,3.15) {$y=3$};
\end{axis}
\end{tikzpicture}
\end{center}

На промежутке
\[
[-2;1]
\]
график горизонтален.

Поэтому прямая
\[
y=m
\]
имеет с графиком \textbf{бесконечно много общих точек} при
\[
\boxed{m=3}.
\]

Множество значений:
\[
E(f)=[3;+\infty).
\]

% -------------------------------------------------

\lessonsection{5. Точки, в которых меняется формула}

Точки смены формулы возникают там, где некоторое подмодульное выражение
становится равным нулю.

Например,
\[
g(x)=|x^2-4|.
\]

Решаем
\[
x^2-4=0,
\]
откуда
\[
x=-2,\qquad x=2.
\]

Знаки выражения:
\[
x^2-4
\begin{cases}
>0, & x<-2,\\
<0, & -2<x<2,\\
>0, & x>2.
\end{cases}
\]

Следовательно,
\[
\boxed{
g(x)=
\begin{cases}
x^2-4, & x\le -2,\\[1mm]
4-x^2, & -2<x<2,\\[1mm]
x^2-4, & x\ge 2.
\end{cases}}
\]

Точки смены формулы по оси \(x\):
\[
\boxed{x=-2,\qquad x=2}.
\]

Соответствующие точки графика:
\[
(-2;0),\qquad (2;0).
\]

\textbf{Не путайте:}
значения \(x=-2\) и \(x=2\) --- это абсциссы точек смены формулы;
координаты точек графика записываются двумя числами.

% -------------------------------------------------

\lessonsection{6. Особый случай: квадратное выражение не меняет знак}

Рассмотрим
\[
P(x)=x^2-3x+6.
\]

Дискриминант:
\[
D=(-3)^2-4\cdot1\cdot6=9-24=-15<0.
\]

Коэффициент при \(x^2\) положителен, поэтому
\[
x^2-3x+6>0
\qquad\text{при всех }x\in\R.
\]

Следовательно,
\[
|x^2-3x+6|=x^2-3x+6.
\]

\important{Отсутствие действительных корней не означает отсутствие решений
неравенства.}

Если квадратный трёхчлен не имеет действительных корней, его знак
определяется знаком старшего коэффициента.

% -------------------------------------------------

\lessonsection{7. Алгебраическая дробь с модулем: сначала ОДЗ}

Рассмотрим функцию
\[
q(x)=
\frac{2-|x-1|}
{(x-1)^2-2|x-1|}.
\]

Поскольку
\[
(x-1)^2=|x-1|^2,
\]
знаменатель раскладывается:
\[
(x-1)^2-2|x-1|
=
|x-1|\bigl(|x-1|-2\bigr).
\]

Поэтому исходная функция определена только при
\[
|x-1|\ne0
\qquad\text{и}\qquad
|x-1|\ne2.
\]

Получаем ограничения:
\[
x\ne1,
\]
\[
x-1\ne2,\qquad x-1\ne-2,
\]
то есть
\[
\boxed{x\ne-1,\quad x\ne1,\quad x\ne3}.
\]

Теперь числитель:
\[
2-|x-1|=
-\bigl(|x-1|-2\bigr).
\]

При допустимых \(x\) можно сократить общий множитель:
\[
q(x)=
\frac{-\bigl(|x-1|-2\bigr)}
{|x-1|\bigl(|x-1|-2\bigr)}
=
-\frac{1}{|x-1|}.
\]

Но после сокращения \important{обязательно сохраняются ограничения исходной дроби}:
\[
\boxed{
q(x)=-\frac1{|x-1|},
\qquad
x\ne-1,1,3.}
\]

Это принципиально важно:
точки \(x=-1\) и \(x=3\) нельзя вернуть в область определения после
сокращения.

% -------------------------------------------------

\lessonsection{8. Кусочная формула дробно-рациональной функции}

Рассмотрим знак \(x-1\).

При
\[
x>1
\]
имеем
\[
|x-1|=x-1,
\]
поэтому
\[
q(x)=-\frac1{x-1}.
\]

При
\[
x<1
\]
имеем
\[
|x-1|=-(x-1),
\]
поэтому
\[
q(x)
=
-\frac1{-(x-1)}
=
\frac1{x-1}.
\]

С учётом исходного ОДЗ:
\[
\boxed{
q(x)=
\begin{cases}
\dfrac1{x-1},
&
x<1,\quad x\ne-1,
\\[3mm]
-\dfrac1{x-1},
&
x>1,\quad x\ne3.
\end{cases}}
\]

% -------------------------------------------------

\lessonsection{9. Асимптоты гиперболы}

Для функции
\[
y=\frac{k_1x+k_2}{k_3x+k_4},
\qquad k_3\ne0,
\]
вертикальная асимптота находится из знаменателя:
\[
k_3x+k_4=0,
\]
то есть
\[
\boxed{x=-\frac{k_4}{k_3}}.
\]

Если степени числителя и знаменателя одинаковы, горизонтальная асимптота:
\[
\boxed{y=\frac{k_1}{k_3}}.
\]

Для
\[
-\frac1{x-1}
\qquad\text{и}\qquad
\frac1{x-1}
\]
получаем:
\[
\boxed{x=1}
\]
--- вертикальная асимптота,
и
\[
\boxed{y=0}
\]
--- горизонтальная асимптота.

\textbf{Важно.}
Асимптота \(x=1\) не является обычной ``выколотой точкой'':
при приближении к \(x=1\) значение функции неограниченно убывает.

Точки \(x=-1\) и \(x=3\), напротив, являются устранимыми разрывами,
возникшими из-за сокращения.

В обеих точках сокращённая формула дала бы
\[
-\frac12,
\]
следовательно, на графике отсутствуют точки
\[
\boxed{\left(-1;-\frac12\right)}
\qquad\text{и}\qquad
\boxed{\left(3;-\frac12\right)}.
\]

\begin{center}
\begin{tikzpicture}
\begin{axis}[
    width=11cm,
    height=7.2cm,
    axis lines=middle,
    unit vector ratio=1 1,
    xmin=-4,
    xmax=6,
    ymin=-4,
    ymax=1.2,
    xlabel={$x$},
    ylabel={$q$},
    samples=300,
    clip=false,
    grid=both,
    major grid style={softgray}
]

% Вертикальная асимптота
\addplot[
    color=muted,
    dashed,
    thick
]
coordinates {(1,-4) (1,1.1)};

% Левая ветвь
\addplot[
    color=accent,
    very thick,
    domain=-4:0.94
]
{1/(x-1)};

% Правая ветвь
\addplot[
    color=accent,
    very thick,
    domain=1.06:6
]
{-1/(x-1)};

% Выколотые точки
\addplot[
    only marks,
    mark=o,
    mark size=3.2pt,
    mark options={line width=1.2pt,fill=white},
    color=darkaccent
]
coordinates {
    (-1,-0.5)
    (3,-0.5)
};

\node[anchor=south west,text=muted]
at (axis cs:1.08,0.55) {$x=1$};

\end{axis}
\end{tikzpicture}
\end{center}

Так как
\[
q(x)=-\frac1{|x-1|}<0,
\]
функция принимает только отрицательные значения.

Без выколотых точек множество значений было бы
\[
(-\infty;0).
\]

Однако значение
\[
-\frac12
\]
достигается только при
\[
|x-1|=2,
\]
то есть при \(x=-1\) или \(x=3\), а обе эти точки запрещены.

Поэтому
\[
\boxed{
E(q)=
\left(-\infty;-\frac12\right)
\cup
\left(-\frac12;0\right).}
\]

% -------------------------------------------------

\lessonsection{10. Построение гиперболы: рабочий алгоритм}

Если требуется построить дробно-рациональную функцию, действуйте так:

\begin{enumerate}
    \item Запишите \important{ОДЗ исходной функции}.
    \item Выполните допустимые алгебраические преобразования.
    \item Если сократили множитель, \important{не удаляйте исходные ограничения}.
    \item Раскройте модули и получите кусочную формулу.
    \item Найдите вертикальные и горизонтальные асимптоты.
    \item Найдите точки пересечения с осями, если они существуют и входят
    в рассматриваемый промежуток.
    \item При необходимости составьте небольшую таблицу значений.
    \item Постройте каждую ветвь только на разрешённом промежутке.
    \item Отметьте выколотые точки.
\end{enumerate}

\textbf{Для эскиза} одной корректно выбранной точки на каждой ветви
часто достаточно, если асимптоты уже известны.

% -------------------------------------------------

\lessonsection{11. Типичные ошибки}

\begin{enumerate}
    \item \textbf{Раскрывать модуль без проверки знака.}

    Нельзя автоматически писать
    \[
    |A(x)|=A(x).
    \]

    \item \textbf{Путать отсутствие корней с отсутствием решений.}

    Из \(D<0\) следует отсутствие действительных корней квадратного
    трёхчлена, а не отсутствие решений любого неравенства с ним.

    \item \textbf{Забывать внешний минус.}

    Например,
    \[
    -|3x+6|
    \]
    при \(3x+6<0\) превращается в
    \[
    -\bigl(-(3x+6)\bigr)=3x+6.
    \]

    \item \textbf{Строить ветвь вне её области существования.}

    Если формула получена для \(x>1\), точки с \(x\le1\)
    использовать для этой ветви нельзя.

    \item \textbf{Терять запрещённые точки после сокращения дроби.}

    Сокращение преобразует формулу, но не расширяет ОДЗ исходной функции.

    \item \textbf{Путать множество значений и область определения.}

    Область определения --- допустимые \(x\).

    Множество значений --- достигаемые \(y\).

    \item \textbf{Путать точку смены формулы с её абсциссой.}

    \(x=2\) --- значение аргумента.

    \((2;0)\) --- точка координатной плоскости.
\end{enumerate}

% =================================================
% ТРЕНИРОВОЧНЫЙ БЛОК
% =================================================
\clearpage

\begin{center}
{\LARGE\bfseries\color{darkaccent}Тренировочный блок}

\vspace{0.3em}

{\color{muted}10 упражнений по нарастающей сложности}
\end{center}

\vspace{0.6em}

\textbf{1. Средний уровень.}
Запишите функцию
\[
y=3-|x|
\]
в кусочном виде и найдите её множество значений.

\textbf{2. Средний уровень.}
Запишите
\[
f(x)=|x-2|
\]
в кусочном виде. Укажите точку смены формулы и множество значений.

\textbf{3. Средний уровень.}
Для функции
\[
g(x)=|x-1|+|x+2|
\]
запишите кусочную формулу и укажите промежуток, на котором функция постоянна.

\vspace{0.5em}

\textbf{4. Выше среднего.}
Запишите в кусочном виде функцию
\[
h(x)=|2x+4|-1.
\]
Укажите точку смены формулы и множество значений.

\textbf{5. Выше среднего.}
Для функции
\[
p(x)=|x^2-4|
\]
запишите кусочную формулу и укажите все точки смены формулы.

\textbf{6. Выше среднего.}
Дана функция
\[
F(x)=|x-4|+|x+1|.
\]
Найдите значение \(m\), при котором графики
\[
y=F(x)
\qquad\text{и}\qquad
y=m
\]
имеют бесконечно много общих точек.

\textbf{7. Выше среднего.}
Упростите функцию
\[
r(x)=
\frac{4-|x|}
{x^2-4|x|}
\]
с обязательным сохранением ОДЗ. Укажите вертикальную и горизонтальную
асимптоты, а также координаты выколотых точек.

\textbf{8. Выше среднего.}
Решите уравнение
\[
|x-3|+|x+1|=6.
\]

\textbf{9. Выше среднего.}
Запишите в кусочном виде функцию
\[
G(x)=|x^2-1|.
\]
Укажите точки смены формулы и найдите \(G(0)\).

\vspace{0.5em}

\textbf{10. Сложный уровень.}
Исследуйте функцию
\[
q(x)=
\frac{2-|x-1|}
{(x-1)^2-2|x-1|}.
\]

Требуется:
\begin{enumerate}[label=\alph*)]
    \item записать ОДЗ исходной функции;
    \item упростить формулу без потери ограничений;
    \item записать функцию в кусочном виде;
    \item найти асимптоты;
    \item указать выколотые точки;
    \item найти множество значений.
\end{enumerate}

% =================================================
% ОТВЕТЫ
% =================================================
\clearpage

\begin{center}
{\LARGE\bfseries\color{darkaccent}Ответы к тренировочному блоку}
\end{center}

\vspace{0.8em}

\renewcommand{\arraystretch}{1.35}

\begin{table}[ht]
\centering
\caption{Краткие ответы}
\begin{tabularx}{\linewidth}{@{}cX@{}}
\toprule
\textbf{№} & \textbf{Ответ} \\
\midrule

1 &
\(
y=
\begin{cases}
3+x,&x<0,\\
3-x,&x\ge0;
\end{cases}
\quad
E(y)=(-\infty;3].
\)
\\

2 &
\(
f(x)=
\begin{cases}
2-x,&x<2,\\
x-2,&x\ge2;
\end{cases}
\)
точка смены: \(x=2\);
\(
E(f)=[0;+\infty).
\)
\\

3 &
\(
g(x)=
\begin{cases}
-2x-1,&x<-2,\\
3,&-2\le x<1,\\
2x+1,&x\ge1.
\end{cases}
\)
Функция постоянна на \([-2;1]\).
\\

4 &
\(
h(x)=
\begin{cases}
-2x-5,&x<-2,\\
2x+3,&x\ge-2.
\end{cases}
\)
Точка смены: \(x=-2\);
\(
E(h)=[-1;+\infty).
\)
\\

5 &
\(
p(x)=
\begin{cases}
x^2-4,&x\le-2,\\
4-x^2,&-2<x<2,\\
x^2-4,&x\ge2.
\end{cases}
\)
Точки смены: \(x=-2,\ 2\).
\\

6 &
\(
m=5.
\)
Горизонтальный участок расположен при
\(
x\in[-1;4].
\)
\\

7 &
\(
r(x)=-\dfrac1{|x|},
\quad
x\ne-4,0,4.
\)
Асимптоты:
\(
x=0,\ y=0.
\)
Выколотые точки:
\(
\left(-4;-\frac14\right),
\left(4;-\frac14\right).
\)
\\

8 &
\(
x=-2,\ 4.
\)
\\

9 &
\(
G(x)=
\begin{cases}
x^2-1,&x\le-1,\\
1-x^2,&-1<x<1,\\
x^2-1,&x\ge1.
\end{cases}
\)
Точки смены:
\(
x=-1,\ 1;
\quad
G(0)=1.
\)
\\

10 &
\(
D(q)=\R\setminus\{-1,1,3\}.
\)

\(
q(x)=-\dfrac1{|x-1|}.
\)

\(
q(x)=
\begin{cases}
\dfrac1{x-1},&x<1,\ x\ne-1,\\[2mm]
-\dfrac1{x-1},&x>1,\ x\ne3.
\end{cases}
\)

Асимптоты:
\(
x=1,\ y=0.
\)

Выколотые точки:
\(
\left(-1;-\frac12\right),
\left(3;-\frac12\right).
\)

\(
E(q)=
\left(-\infty;-\frac12\right)
\cup
\left(-\frac12;0\right).
\)
\\

\bottomrule
\end{tabularx}
\end{table}

\end{document}